\documentclass[10pt,a4paper]{amsart}
\usepackage{amsmath,amssymb,amsthm,graphicx}
\usepackage{verbatim}
\newenvironment{reference}{\comment}{\endcomment}
%\newenvironment{reference}{}{}
\newenvironment{slogan}{\comment}{\endcomment}
\newenvironment{history}{\comment}{\endcomment}
% For commutative diagrams we use Xy-pic
\usepackage[all]{xy}

% We use 2cell for 2-commutative diagrams.
\xyoption{2cell}
% Theorem environments.
%
\theoremstyle{plain}
\newtheorem{theorem}[subsection]{Theorem}
\newtheorem{proposition}[subsection]{Proposition}
\newtheorem{lemma}[subsection]{Lemma}

\theoremstyle{definition}
\newtheorem{definition}[subsection]{Definition}
\newtheorem{example}[subsection]{Example}
\newtheorem{exercise}[subsection]{Exercise}
\newtheorem{situation}[subsection]{Situation}

\theoremstyle{remark}
\newtheorem{remark}[subsection]{Remark}
\newtheorem{remarks}[subsection]{Remarks}

\numberwithin{equation}{subsection}

% Macros
%
\def\lim{\mathop{\mathrm{lim}}\nolimits}
\def\colim{\mathop{\mathrm{colim}}\nolimits}
\def\Spec{\mathop{\mathrm{Spec}}}
\def\Hom{\mathop{\mathrm{Hom}}\nolimits}
\def\Ext{\mathop{\mathrm{Ext}}\nolimits}
\def\SheafHom{\mathop{\mathcal{H}\!\mathit{om}}\nolimits}
\def\SheafExt{\mathop{\mathcal{E}\!\mathit{xt}}\nolimits}
\def\Sch{\mathit{Sch}}
\def\Mor{\mathop{\mathrm{Mor}}\nolimits}
\def\Ob{\mathop{\mathrm{Ob}}\nolimits}
\def\Sh{\mathop{\mathit{Sh}}\nolimits}
\def\NL{\mathop{N\!L}\nolimits}
\def\CH{\mathop{\mathrm{CH}}\nolimits}
\def\proetale{{pro\text{-}\acute{e}tale}}
\def\etale{{\acute{e}tale}}
\def\QCoh{\mathit{QCoh}}
\def\Ker{\mathop{\mathrm{Ker}}}
\def\Im{\mathop{\mathrm{Im}}}
\def\Coker{\mathop{\mathrm{Coker}}}
\def\Coim{\mathop{\mathrm{Coim}}}

% Boxtimes
%
\DeclareMathSymbol{\boxtimes}{\mathbin}{AMSa}{"02}

%
% Macros for moduli stacks/spaces
%
\def\QCohstack{\mathcal{QC}\!\mathit{oh}}
\def\Cohstack{\mathcal{C}\!\mathit{oh}}
\def\Spacesstack{\mathcal{S}\!\mathit{paces}}
\def\Quotfunctor{\mathrm{Quot}}
\def\Hilbfunctor{\mathrm{Hilb}}
\def\Curvesstack{\mathcal{C}\!\mathit{urves}}
\def\Polarizedstack{\mathcal{P}\!\mathit{olarized}}
\def\Complexesstack{\mathcal{C}\!\mathit{omplexes}}
% \Pic is the operator that assigns to X its picard group, usage \Pic(X)
% \Picardstack_{X/B} denotes the Picard stack of X over B
% \Picardfunctor_{X/B} denotes the Picard functor of X over B
\def\Pic{\mathop{\mathrm{Pic}}\nolimits}
\def\Picardstack{\mathcal{P}\!\mathit{ic}}
\def\Picardfunctor{\mathrm{Pic}}
\def\Deformationcategory{\mathcal{D}\!\mathit{ef}}
\usepackage{cleveref}
\crefname{originallemma}{original lemma}{original lemmas}
\Crefname{originallemma}{Original lemma}{Original lemmas}
\makeatletter
\newlabel{lemma-dualizing-localize}{{lemma-dualizing-localize}{source}{}{sourcepointer-1}{}}
\newlabel{lemma-dualizing-localize@cref}{{[originallemma][0][]lemma-dualizing-localize}{[1][0][]source}{}{}{}}
\makeatother
\crefname{originallemma}{original lemma}{original lemmas}
\Crefname{originallemma}{Original lemma}{Original lemmas}
\makeatletter
\newlabel{lemma-normalized-quotient}{{lemma-normalized-quotient}{source}{}{sourcepointer-2}{}}
\newlabel{lemma-normalized-quotient@cref}{{[originallemma][0][]lemma-normalized-quotient}{[1][0][]source}{}{}{}}
\makeatother
\crefname{originallemma}{original lemma}{original lemmas}
\Crefname{originallemma}{Original lemma}{Original lemmas}
\makeatletter
\newlabel{lemma-sitting-in-degrees}{{lemma-sitting-in-degrees}{source}{}{sourcepointer-3}{}}
\newlabel{lemma-sitting-in-degrees@cref}{{[originallemma][0][]lemma-sitting-in-degrees}{[1][0][]source}{}{}{}}
\makeatother
\crefname{originalreference}{original reference}{original references}
\Crefname{originalreference}{Original reference}{Original references}
\makeatletter
\newlabel{more-algebra-lemma-base-change-RHom}{{more-algebra-lemma-base-change-RHom}{source}{}{sourcepointer-4}{}}
\newlabel{more-algebra-lemma-base-change-RHom@cref}{{[originalreference][0][]more-algebra-lemma-base-change-RHom}{[1][0][]source}{}{}{}}
\makeatother
\title{The degree in which a normalized dualizing complex is nonzero at a minimal prime}
\author{The Stacks Project Authors}
\date{stacks-project:dualizing.tex:local-snapshot-20261012; GNU FDL 1.2 or later}
\begin{document}
\maketitle
\noindent\textit{Editorial note.} All delivered mathematical text is the author's own, copied from the pinned Stacks Project chapter source sources/stacks-dualizing/pinned-original-source.tex; no statement or proof step was rewritten, reordered or removed.\par\medskip
\noindent\textit{Editorial note.} Preamble provenance: the Stacks Project keeps one single shared preamble for every chapter; this card takes it directly from the staged sources/stacks-dualizing/preamble.tex as three verbatim spans (lines 8-12, 14-18 and 154-217), so no derived file was written and nothing was added or redefined. The wrapper class and the base packages are supplied by the assembler.\par\medskip
\noindent\textit{Editorial note.} Bibliography provenance: neither the counted statement nor its proof contains a citation (this lane rescans the whole span before writing the card), so the card carries no bibliography and no entry was reconstructed.\par\medskip
\noindent\textit{Editorial note.} Selection provenance: this card was opened under the file-based control contract (CONTROL.md), with the mechanical pre-card gate run on its statement and proof interval before the card was written.\par\medskip
\noindent\textit{Allowed standard prerequisites (editorial).}\begin{itemize}
\item The quotient by a power of a prime ideal of dimension zero inherits a normalized dualizing complex. Invoked inside the counted proof (source lines 3426). Replaces the ring by a quotient in which the ideal has become zero, without changing the localized complex.
\item The degree confinement for the cohomology of a normalized dualizing complex. Invoked inside the counted proof (source lines 3428-3429). Shows that only the degree equal to minus the dimension can be nonzero.
\item The localization of a dualizing complex at a prime is dualizing over the local ring. Invoked inside the counted proof (source lines 3430-3432). Shows the remaining cohomology module is nonzero.
\end{itemize}\medskip
\ifdefined\hypertarget\hypertarget{sourcepointer-1}{}\fi
\noindent\textit{Source reference.} Original lemma lemma-dualizing-localize: sources/stacks-dualizing/pinned-original-source.tex, source lines 2950--2950. Navigation pointer only: the counted text refers to the Stacks item lemma-dualizing-localize, which is an existing labeled result of this same chapter and is not part of this card. The pointer carries the target's own label as its number and its definition line as its span; its content is not reproduced here.\par\medskip
\ifdefined\hypertarget\hypertarget{sourcepointer-2}{}\fi
\noindent\textit{Source reference.} Original lemma lemma-normalized-quotient: sources/stacks-dualizing/pinned-original-source.tex, source lines 3160--3160. Navigation pointer only: the counted text refers to the Stacks item lemma-normalized-quotient, which is an existing labeled result of this same chapter and is not part of this card. The pointer carries the target's own label as its number and its definition line as its span; its content is not reproduced here.\par\medskip
\ifdefined\hypertarget\hypertarget{sourcepointer-3}{}\fi
\noindent\textit{Source reference.} Original lemma lemma-sitting-in-degrees: sources/stacks-dualizing/pinned-original-source.tex, source lines 3232--3232. Navigation pointer only: the counted text refers to the Stacks item lemma-sitting-in-degrees, which is an existing labeled result of this same chapter and is not part of this card. The pointer carries the target's own label as its number and its definition line as its span; its content is not reproduced here.\par\medskip
\ifdefined\hypertarget\hypertarget{sourcepointer-4}{}\fi
\noindent\textit{Source reference.} Original reference more-algebra-lemma-base-change-RHom: sources/stacks-dualizing/pinned-original-source.tex, source lines 2719--2719. Navigation pointer only: the counted text refers to the Stacks item more-algebra-lemma-base-change-RHom, which belongs to another chapter of the project and is not part of this card. The pointer carries the target's own label as its number because this chapter's staged snapshot does not contain the other chapter's tag table; its content is not reproduced here.\par\medskip
\begin{lemma}
\label{lemma-nonvanishing-generically-local}
Let $(A, \mathfrak m, \kappa)$ be a Noetherian local ring with
normalized dualizing complex $\omega_A^\bullet$.
Let $\mathfrak p$ be a minimal prime of $A$ with
$\dim(A/\mathfrak p) = e$. Then
$H^i(\omega_A^\bullet)_\mathfrak p$ is nonzero
if and only if $i = -e$.
\end{lemma}

\begin{proof}
Since $A_\mathfrak p$ has dimension zero, there exists an integer
$n > 0$ such that $\mathfrak p^nA_\mathfrak p$ is zero.
Set $B = A/\mathfrak p^n$ and
$\omega_B^\bullet = R\Hom_A(B, \omega_A^\bullet)$.
Since $B_\mathfrak p = A_\mathfrak p$ we see that
$$
(\omega_B^\bullet)_\mathfrak p =
R\Hom_A(B, \omega_A^\bullet) \otimes_A^\mathbf{L} A_\mathfrak p =
R\Hom_{A_\mathfrak p}(B_\mathfrak p, (\omega_A^\bullet)_\mathfrak p) =
(\omega_A^\bullet)_\mathfrak p
$$
The second equality holds by
More on Algebra, Lemma \ref{more-algebra-lemma-base-change-RHom}.
By Lemma \ref{lemma-normalized-quotient} we may replace $A$ by $B$.
After doing so, we see that $\dim(A) = e$. Then we see that
$H^i(\omega_A^\bullet)_\mathfrak p$ can only be nonzero if $i = -e$
by Lemma \ref{lemma-sitting-in-degrees} parts (1) and (2).
On the other hand, since $(\omega_A^\bullet)_\mathfrak p$
is a dualizing complex for the nonzero ring $A_\mathfrak p$
(Lemma \ref{lemma-dualizing-localize})
we see that the remaining module has to be nonzero.
\end{proof}

\end{document}
